\chapter{Comment les obligations se négocient}\label{ch:m2:how-bonds-trade}

Une grande entreprise peut avoir cinquante obligations en circulation, et son action
n'occupe qu'une seule ligne à la cote d'une bourse. Chaque obligation est un titre distinct, doté de son propre coupon et de son échéance, et la plupart
d'entre elles ne sont pas négociées la plupart des jours. Il n'existe pas de carnet d'ordres à consulter~; un
investisseur qui veut vendre demande des prix aux dealers, et la réponse dépend de
qui est sollicité, combien ils sont et ce qu'ils apprennent du seul fait d'être sollicités. Jusqu'en 2002, même les
prix des transactions passées étaient confidentiels. Depuis, les dealers américains ont déclaré
chaque transaction sur obligations d'entreprise vers un flux public, et depuis le 15 mai 2023 ils
signalent les \omterm{def:m2:how-bonds-trade:a2a}{transactions de portefeuille}~: pendant leur première semaine, elles représentaient déjà
5.2~\% du notionnel échangé. Ce chapitre explique la \omterm{def:m2:how-bonds-trade:rfq}{demande de cotation} (RFQ), la \omterm{def:m2:how-bonds-trade:trace}{déclaration
des transactions}, les nouvelles façons dont les obligations s'échangent en bloc ou directement entre elles, les
fonds indiciels cotés (ETF) sur obligations qui se superposent au marché, et comment les dealers
fixent le prix de ce qui se négocie rarement.

\section{Demande de cotation (RFQ)}

\begin{definition}[Demande de cotation, axe]\label{def:m2:how-bonds-trade:rfq}
Une \emph{demande de cotation}\index{demande de cotation} (RFQ) est une requête du client,
généralement envoyée par voie électronique à plusieurs dealers à la fois, pour obtenir un prix ferme d'achat ou
de vente d'un montant donné d'un titre~; le client traite avec la meilleure réponse ou
aucune. Un \emph{axe}\index{axe} est l'intérêt affiché d'un dealer à acheter ou
vendre une obligation donnée, généralement en raison de son inventaire.
\end{definition}

La RFQ est une enchère que le client organise pour chaque transaction. Interroger davantage de dealers apporte
une concurrence accrue et un meilleur prix. Mais chaque dealer interrogé apprend qu'un client
veut vendre cette obligation tout de suite, et un dealer perdant peut exploiter cette information, en vendant
l'obligation ou ses couvertures avant le gagnant, ou simplement la déprécier~; le client
en supporte le coût, dans le prix de cette transaction ou de la suivante. Les \omterm{def:m2:how-bonds-trade:rfq}{axes} indiquent au
client quels dealers veulent la contrepartie et sont donc susceptibles d'offrir un bon prix.

\begin{proposition}[Combien de dealers interroger]\label{prop:m2:how-bonds-trade:n}
Supposons que le dealer $i$ propose $v - m + \sigma Z_i$, avec $v$ la valeur de l'obligation, $m$ une
marge moyenne, et $Z_i$ des termes privés indépendants de loi normale centrée réduite (inventaire, \omterm{def:m2:how-bonds-trade:rfq}{axes}, anticipations),
et que chaque dealer perdant coûte au client $\ell$ par fuite d'information. La vente
à la meilleure des $n$ offres d'achat coûte, en espérance,
\[
C(n) \;=\; m - \sigma\,e_n + \ell\,(n-1),
\]
où $e_n$ est l'espérance du maximum de $n$ normales standard ($e_1 = 0$, $e_2 = 0.564$,
$e_3 = 0.846$, $e_4 = 1.029$). Comme $e_n$ croît de plus en plus lentement, $C$ admet un
minimum pour un $n$ fini, qui diminue à mesure que la fuite $\ell$ augmente.
\end{proposition}

\begin{proof}
La meilleure offre vaut $v - m + \sigma\max_i Z_i$, dont la décote espérée par rapport à $v$ est
$m - \sigma e_n$~; les $n-1$ perdants ajoutent chacun $\ell$. Les accroissements $e_{n+1} - e_n$
décroissent vers zéro, alors que la fuite ajoute $\ell$ pour chaque dealer, donc $C(n+1) - C(n) =
\ell - \sigma(e_{n+1} - e_n)$ devient positive.
\end{proof}

\begin{omfigure}
\begin{tikzpicture}
\begin{axis}[width=11cm,height=4.8cm,
  xlabel={nombre de dealers interrogés, $n$},ylabel={coût espéré (pb)},
  tick label style={font=\small},label style={font=\small},
  xmin=0.7,xmax=8.3,ymin=5.5,ymax=14,grid=major,grid style={gray!25},
  legend columns=3,legend style={font=\footnotesize,at={(0.5,-0.3)},anchor=north,draw=none,fill=none,/tikz/every even column/.append style={column sep=8pt}},
  table/col sep=comma]
\addplot[omDef,very thick,mark=*,mark size=1.4pt] table[x=n,y=l05]{figdata/markets-2/22-how-bonds-trade/cost.csv};
\addplot[omThm,very thick,mark=*,mark size=1.4pt] table[x=n,y=l10]{figdata/markets-2/22-how-bonds-trade/cost.csv};
\addplot[omProp,very thick,mark=*,mark size=1.4pt] table[x=n,y=l15]{figdata/markets-2/22-how-bonds-trade/cost.csv};
\legend{fuite 0.5 pb,fuite 1.0 pb,fuite 1.5 pb}
\end{axis}
\end{tikzpicture}

\omcaption{Coût espéré, pour le client, de la vente d'une obligation par RFQ, en points de
base du prix, selon le nombre de dealers interrogés, avec une marge moyenne de 10 points de base, un
écart de termes privés de 5 et trois niveaux de fuite par dealer perdant. Le
nombre optimal de dealers est de 6, 3 et 2. Illustratif~; données~: le
tutoriel du chapitre.}\label{fig:m2:how-bonds-trade:cost}
\end{omfigure}

\begin{example}[Trois dealers]\label{ex:m2:how-bonds-trade:three}
Avec une marge de 10 points de base, un écart de termes privés de 5 et une fuite de 1 point
de base par dealer perdant, interroger un seul dealer coûte 10.0 points de base, deux 8.18, trois 7.77,
quatre 7.85 et six 8.66. Trois est le meilleur choix~; pour une vente de 25~millions de USD, cela coûte environ
19\,400~USD, contre 25\,000~USD avec un seul dealer.
\end{example}

\begin{omfigure}
\begin{tikzpicture}
\begin{axis}[width=11cm,height=4.8cm,
  xlabel={décote de l'offre gagnante par rapport à la valeur (pb)},ylabel={densité},
  tick label style={font=\small},label style={font=\small},
  xmin=-10,xmax=26,ymin=0,ymax=0.14,grid=major,grid style={gray!25},
  yticklabel style={/pgf/number format/.cd,fixed,precision=2},
  legend columns=3,legend style={font=\footnotesize,at={(0.5,-0.3)},anchor=north,draw=none,fill=none,/tikz/every even column/.append style={column sep=8pt}},
  table/col sep=comma]
\addplot[omDef,thick,const plot mark mid] table[x=x,y=n1]{figdata/markets-2/22-how-bonds-trade/hist.csv};
\addplot[omThm,very thick,const plot mark mid] table[x=x,y=n3]{figdata/markets-2/22-how-bonds-trade/hist.csv};
\addplot[omProp,thick,dashed,const plot mark mid] table[x=x,y=n6]{figdata/markets-2/22-how-bonds-trade/hist.csv};
\legend{1 dealer,3 dealers,6 dealers}
\end{axis}
\end{tikzpicture}

\omcaption{Distribution simulée de la décote de l'offre gagnante par rapport à la valeur de l'obligation,
avant fuite, lorsqu'un seul, trois ou six dealers sont interrogés. Davantage de dealers déplacent la
distribution vers de meilleurs prix et la resserrent, avec des gains décroissants.
Illustratif~; données~: le tutoriel du chapitre.}\label{fig:m2:how-bonds-trade:hist}
\end{omfigure}

\section{Déclaration des transactions et transparence}

\begin{definition}[Déclaration des transactions, plafond de publication]\label{def:m2:how-bonds-trade:trace}
La \emph{déclaration des transactions}\index{déclaration des transactions} est l'obligation pour un dealer de déclarer
chaque transaction, avec son prix, sa taille et son horodatage, au système d'un régulateur qui
la publie. Un \emph{plafond de publication}\index{plafond de publication} est une taille au-delà de laquelle
le compte rendu publié indique seulement que la transaction était supérieure au plafond.
\end{definition}

Le système TRACE de la FINRA collecte les transactions sur obligations d'entreprises américaines depuis juillet 2002~;
les dealers doivent déclarer dans les quinze minutes, et les comptes rendus sont publiés dès leur
réception. La publication s'est faite par étapes sur trois ans, les obligations les plus importantes et les mieux notées
en premier, et les transactions de grande taille ont été publiées avec des tailles plafonnées~: montants exacts
jusqu'à 5~millions de dollars pour les obligations investment grade et 1~million de dollars pour les obligations high yield, «~5MM$+$~»
et «~1MM$+$~» au-delà. Le plafond protège un dealer qui a acheté un gros
bloc contre le front-running pendant qu'il le revend~; le prix publié indique toujours au marché
où se négocient les obligations. L'Union européenne ne disposait pas de système consolidé de publication de ses
transactions obligataires~; en juillet 2025, son régulateur des marchés financiers a retenu le premier fournisseur d'un tel système.

\begin{dated}{2026-09}{Règles de transparence}\label{dat:m2:how-bonds-trade:transparency}
États-Unis~: déclaration TRACE dans les 15~minutes pour les obligations d'entreprises~;
modificateur de \omterm{def:m2:how-bonds-trade:a2a}{transaction de portefeuille} depuis le 15~mai 2023. Union européenne~: l'ESMA a retenu Ediphy
(fairCT) le 3~juillet 2025 comme premier fournisseur de système consolidé de publication pour les obligations, habilité à exploiter le système
pendant cinq ans après agrément.
\end{dated}

\section{Négociation all-to-all et transactions de portefeuille}

\begin{definition}[Négociation all-to-all, transaction de portefeuille]\label{def:m2:how-bonds-trade:a2a}
La \emph{négociation all-to-all}\index{négociation all-to-all} permet à tout participant d'une
plateforme, investisseur ou dealer, de répondre à la demande ou à l'ordre de n'importe quel autre, au lieu
que les investisseurs n'interrogent que les dealers. Une \emph{transaction de portefeuille}\index{transaction de portefeuille} est l'achat ou la vente
d'un panier d'obligations à un prix unique convenu pour l'ensemble du panier~; TRACE
signale une transaction comme telle lorsqu'elle se fait entre deux parties, pour au moins dix
obligations différentes, à un même prix.
\end{definition}

La \omterm{def:m2:how-bonds-trade:a2a}{négociation all-to-all} élargit l'ensemble des acheteurs possibles~: une société de gestion qui vend une
obligation peut trouver une autre société de gestion qui la veut, sans dealer intermédiaire. La négociation
de portefeuille fait l'inverse de la RFQ pour les obligations isolées~: au lieu de demander cent
cotations, le client n'en demande qu'une, pour la liste entière, et le dealer valorise le risque
du panier, qu'il peut couvrir avec des indices de crédit et des ETF obligataires, plutôt que
obligation par obligation. Le prix de chaque obligation dans une \omterm{def:m2:how-bonds-trade:a2a}{transaction de portefeuille} n'est pas un prix négocié
de cette obligation, raison pour laquelle la FINRA a exigé qu'elle soit signalée~: le lecteur du flux ne doit
pas le prendre pour tel.

\begin{omfigure}
\centering
\begin{tikzpicture}[font=\small,
  box/.style={draw,thick,rounded corners=2pt,align=center,minimum height=1cm,minimum width=2.6cm,inner sep=3pt},
  arr/.style={-{Stealth[length=2.5mm]},thick}]
\node[box,draw=omDef,fill=omDef!8] (cl) at (0,0) {client avec une\\liste de 150 obligations};
\node[box,draw=omThm,fill=omThm!8] (dl) at (5.5,0) {dealer de\\transactions de portefeuille};
\node[box,draw=omProp,fill=omProp!8] (etf) at (11,1.1) {ETF obligataire\\(créations, rachats)};
\node[box,draw=gray,fill=gray!8] (cdx) at (11,-1.1) {swaps sur\\indices de crédit};
\draw[arr] (cl) -- node[above,font=\footnotesize,align=center] {un prix\\pour la liste} (dl);
\draw[arr] (dl) -- node[above,sloped,font=\footnotesize] {couverture} (etf);
\draw[arr] (dl) -- node[below,sloped,font=\footnotesize] {couverture} (cdx);
\node[font=\footnotesize,text=omThm,align=center] at (5.5,-1.4) {écoule les obligations\\sur plusieurs jours};
\end{tikzpicture}

\omcaption{Une \omterm{def:m2:how-bonds-trade:a2a}{transaction de portefeuille}. Le client demande un prix unique pour une liste entière~; le dealer
prend le panier, couvre son risque global avec des ETF obligataires et des swaps sur indices, puis vend ou
rachète les obligations individuelles dans les jours qui suivent. Schéma.}\label{fig:m2:how-bonds-trade:pt}
\end{omfigure}

\section{Les ETF obligataires comme couche de liquidité}

Un fonds indiciel coté obligataire (livre~1, chapitre~14) détient des centaines ou des milliers
d'obligations et se négocie en bourse toute la journée. Ses actions se négocient beaucoup plus facilement que
la plupart de ses obligations, et ses participants autorisés créent et rachètent des actions contre
des paniers d'obligations~: lorsque l'ETF se négocie en dessous de la valeur de ses obligations, ils achètent des actions,
les rachètent contre des obligations et vendent les obligations~; lorsqu'il se négocie au-dessus, ils font l'inverse.
L'ETF constitue ainsi une couche liquide au-dessus d'un marché peu liquide~: un investisseur peut modifier
son exposition au crédit en quelques minutes, et le travail lourd de négociation des obligations revient
aux participants, par paniers (\cref{fig:m2:how-bonds-trade:etf}). Dans les marchés sous tension,
le cours de l'ETF peut passer sous la valeur publiée de ses obligations, dont les cours sont
périmés faute d'avoir été négociés, et l'ETF devient le lieu où les prix
se découvrent.

\begin{omfigure}
\centering
\begin{tikzpicture}[font=\small,
  box/.style={draw,thick,rounded corners=2pt,align=center,minimum height=1cm,minimum width=2.6cm,inner sep=3pt},
  arr/.style={{Stealth[length=2.5mm]}-{Stealth[length=2.5mm]},thick}]
\node[box,draw=omDef,fill=omDef!8] (inv) at (0,0) {investisseurs\\(négocient des actions\\toute la journée)};
\node[box,draw=omProp,fill=omProp!8] (ex) at (4.3,0) {bourse~:\\actions d'ETF};
\node[box,draw=omThm,fill=omThm!8] (ap) at (8.6,0) {participants\\autorisés};
\node[box,draw=gray,fill=gray!8] (fund) at (8.6,-2.6) {le fonds\\(créations, rachats)};
\node[box,draw=gray,fill=gray!8] (otc) at (12.9,0) {marché obligataire\\(RFQ, portefeuilles)};
\draw[arr] (inv) -- (ex); \draw[arr] (ex) -- (ap); \draw[arr] (ap) -- node[right,font=\footnotesize,align=left] {actions contre\\paniers} (fund);
\draw[arr] (ap) -- node[above,font=\footnotesize] {obligations} (otc);
\end{tikzpicture}

\omcaption{Un ETF obligataire comme couche de liquidité. Les investisseurs négocient des actions en bourse~;
les participants autorisés échangent des actions contre des paniers d'obligations avec le fonds, et achètent
ou vendent les obligations sur le marché de gré à gré, ce qui maintient le cours de l'action près de la
valeur des obligations. Schéma.}\label{fig:m2:how-bonds-trade:etf}
\end{omfigure}

\section{Tarification algorithmique par les dealers}

\begin{definition}[Prix composite]\label{def:m2:how-bonds-trade:composite}
Un \emph{prix composite}\index{prix composite} d'une obligation est un prix estimé unique
construit à partir de nombreuses observations, transactions récentes, cotations de dealers et prix d'obligations apparentées,
pondérés par leur fraîcheur et leur fiabilité, les valeurs aberrantes étant écartées.
\end{definition}

Un dealer qui répond à des milliers de RFQ par jour ne peut pas tarifer chacune à la main. Son moteur de tarification
part d'un \omterm{def:m2:how-bonds-trade:composite}{prix composite} pour chaque obligation, ancré sur les dernières transactions déclarées à
TRACE et sur les cotations, le déplace avec les autres obligations de l'émetteur, ses CDS et l'indice
pertinent, puis ajoute une marge qui dépend de la taille, du sens, du
client, ainsi que de l'inventaire du dealer et de ses \omterm{def:m2:how-bonds-trade:rfq}{axes}. Le même moteur détermine combien de RFQ traiter et
avec quelle agressivité, l'autre versant de la
\cref{prop:m2:how-bonds-trade:n}.

\begin{example}[Un prix composite]\label{ex:m2:how-bonds-trade:composite}
Des cotations de 99.42, 99.45, 99.40 et 99.47, âgées de 10, 40, 90 et 5~secondes, et une
transaction périmée publiée à 98.20, âgée de 20~secondes~: la valeur aberrante est écartée par la règle de l'écart médian,
et la pondération du reste par une demi-vie d'une minute donne un \omterm{def:m2:how-bonds-trade:composite}{prix composite} de 99.441.
\end{example}

\section{Tutoriel~: un simulateur de RFQ}\label{tut:m2:how-bonds-trade:rfq}

\begin{tutorial}
\textbf{Objectif.} Simuler des enchères RFQ, calculer le coût espéré du client en fonction du nombre de dealers, avec fuite d'information,
trouver le meilleur nombre et construire un \omterm{def:m2:how-bonds-trade:composite}{prix composite}.
\textbf{État final~:} \cref{fig:m2:how-bonds-trade:cost,fig:m2:how-bonds-trade:hist},
\cref{ex:m2:how-bonds-trade:three,ex:m2:how-bonds-trade:composite} et les nombres du
problème du week-end.
\begin{tutsteps}
\item \textbf{L'enchère}~: meilleur cours acheteur espéré, coût avec fuite d'information, le meilleur $n$.
\omcode{code/firm/rfq/firm_rfq.py}{14}{40}{Meilleur cours acheteur espéré, coût pour le client avec fuite d'information, meilleur nombre de dealers, simulation.}\label{lst:m2:how-bonds-trade:rfq}
\item \textbf{Le \omterm{def:m2:how-bonds-trade:composite}{prix composite}}~: valeurs aberrantes écartées, pondérations selon la fraîcheur.
\omcode{code/firm/rfq/firm_rfq.py}{43}{51}{Un prix composite à partir de cotations d'âges différents.}\label{lst:m2:how-bonds-trade:composite}
\item \textbf{Exécuter} \texttt{rfq\_demo.problem()}, \texttt{rfq\_demo.composite\_example()}
et \texttt{fig\_rfq.py}.
\end{tutsteps}
\textbf{Pour aller plus loin.} Faire croître la fuite d'information avec la taille de la transaction et constater
que le meilleur nombre de dealers diminue pour les grandes transactions~; puis donner à un dealer un \omterm{def:m2:how-bonds-trade:rfq}{axe} (un
cours acheteur moyen plus élevé) et déterminer quand interroger ce seul dealer est le meilleur choix.
\end{tutorial}

\section{Construction~: le simulateur de RFQ}\label{bld:m2:how-bonds-trade:rfq}

\begin{build}
\bfield{Rôle} La petite société demande et répond à la fois à des RFQ sur obligations~: en
tant que client, elle choisit combien de dealers interroger~; en tant que dealer, elle fixe ses cotations à partir de \omterm{def:m2:how-bonds-trade:composite}{prix composites}.
\bfield{Interface} \texttt{expected\_max\_normal(n)}~; \texttt{expected\_cost(n, markup,
sigma, leak)}~; \texttt{best\_n(markup, sigma, leak, n\_max)}~;
\texttt{simulate\_rfq(n, markup, sigma, trials, seed)}~; \texttt{composite(quotes,
half\_life, k)}.
\bfield{Règles} Coûts en points de base du prix~; termes privés normaux indépendants~;
fuite d'information linéaire en fonction du nombre de perdants~; \omterm{def:m2:how-bonds-trade:composite}{prix composite} avec retrait des valeurs
aberrantes par écart médian et pondérations exponentielles selon la fraîcheur.
\bfield{Tests d'acceptation} \texttt{code/firm/rfq/tests/}~: $e_1 = 0$ et $e_2 \approx
1/\sqrt\pi$~; sans fuite, il est préférable d'interroger beaucoup de dealers, avec une forte fuite un seul~; la moyenne simulée
coïncide avec la formule~; le \omterm{def:m2:how-bonds-trade:composite}{prix composite} ignore une valeur aberrante.
\bfield{Pour aller plus loin} Des dealers qui apprennent l'historique du client~; termes privés corrélés~;
un \omterm{def:m2:how-bonds-trade:composite}{prix composite} ajusté sur toute la courbe de l'émetteur~; modèles de réponse aux RFQ entraînés sur
les données des plateformes (livre~12).
\end{build}

\begin{omsources}
\item FINRA, TRACE overview~; Regulatory Notice 22-12 on portfolio trades.
\item P.~Asquith, T.~Covert et P.~Pathak, «~The effects of mandatory transparency in
financial market design: evidence from the corporate bond market~», NBER Working
Paper 19417, online appendix.
\item ICE, «~Early observations about the new portfolio trade flag from FINRA TRACE~»,
mai 2023.
\item ESMA, selection of the consolidated tape provider for bonds, juillet 2025.
\end{omsources}

\section{Exercices}

\begin{exercise}[$\star$]\label{exo:m2:how-bonds-trade:1}
Une publication TRACE montre une obligation investment grade négociée «~5MM$+$~» à 99.50.
Déduire ce que l'on sait de la taille de la transaction.
\end{exercise}

\begin{exercise}[$\star$]\label{exo:m2:how-bonds-trade:2}
Avec les paramètres du chapitre et sans fuite d'information, calculer le coût espéré
de la demande adressée à un, deux et quatre dealers.
\end{exercise}

\begin{exercise}[$\star$]\label{exo:m2:how-bonds-trade:3}
Expliquer pourquoi les \omterm{def:m2:how-bonds-trade:a2a}{transactions de portefeuille} doivent être signalées sur le tape.
\end{exercise}

\begin{exercise}[$\star\star$]\label{exo:m2:how-bonds-trade:4}
Donner le meilleur nombre de dealers à interroger lorsque la fuite vaut 0.5, 1.0 et 1.5 points de base
par perdant.
\end{exercise}

\begin{exercise}[$\star\star$]\label{exo:m2:how-bonds-trade:5}
Un ETF obligataire s'échange à 2~\% sous la valeur de ses obligations lors d'une semaine de tensions.
Décrire le rôle d'un participant autorisé et expliquer pourquoi il peut hésiter.
\end{exercise}

\begin{exercise}[$\star\star$]\label{exo:m2:how-bonds-trade:6}
Expliquer pourquoi un dealer peut ne pas répondre du tout à une RFQ.
\end{exercise}

\begin{exercise}[$\star\star\star$]\label{exo:m2:how-bonds-trade:7}
\emph{Programmation.} Avec \texttt{composite}, recalculer \cref{ex:m2:how-bonds-trade:composite}
avec une demi-vie de 10~secondes, puis expliquer le changement.
\end{exercise}

\begin{exercise}[$\star\star\star$]\label{exo:m2:how-bonds-trade:8}
\emph{Repérer l'erreur.} «~Toujours envoyer la RFQ à chaque dealer de la plateforme~: davantage
de concurrence ne peut que donner un meilleur prix.~» Corriger cette affirmation.
\end{exercise}

\section{Problème~: combien de dealers~?}

\begin{problem}[{Problème du week-end --- le coût de la concurrence et d'être repéré}]\label{pb:m2:how-bonds-trade:1}
Une société de gestion doit vendre 25~millions de dollars d'une obligation d'entreprise. Les cours acheteurs des dealers s'établissent en moyenne 10
points de base sous la valeur, avec des termes privés dispersés de 5 points de base~; chaque dealer qui
voit la demande et perd l'enchère coûte en moyenne 1 point de base à la société de gestion.

\medskip\noindent\textbf{Partie I --- L'enchère.}
\begin{enumerate}
\item Donner le coût espéré de la demande adressée à un seul dealer, en points de base et en dollars.
\item Donner les coûts espérés pour deux, trois, quatre et six dealers.
\item Déterminer le meilleur nombre de dealers et son coût.
\item Calculer l'économie réalisée par rapport à six dealers.
\item Expliquer pourquoi le gain apporté par chaque dealer supplémentaire diminue.
\end{enumerate}

\medskip\noindent\textbf{Partie II --- Fuite d'information.}
\begin{enumerate}[resume]
\item Expliquer comment ce que sait un dealer perdant coûte à la société de gestion.
\item Donner le meilleur nombre de dealers si la fuite vaut 0.5 et 1.5 points de base.
\item Expliquer pourquoi la fuite croît avec la taille de la transaction.
\item Expliquer comment les \omterm{def:m2:how-bonds-trade:rfq}{axes} modifient le choix.
\item Expliquer comment la \omterm{def:m2:how-bonds-trade:a2a}{négociation all-to-all} modifierait ce choix.
\end{enumerate}

\medskip\noindent\textbf{Partie III --- Alternatives.}
\begin{enumerate}[resume]
\item Préciser quand le gestionnaire inclurait plutôt l'obligation dans une \omterm{def:m2:how-bonds-trade:a2a}{transaction de portefeuille}.
\item Préciser ce que montre TRACE au sujet de la vente, et à quel moment.
\item Expliquer comment le \omterm{def:m2:how-bonds-trade:trace}{plafond de publication} protège le dealer gagnant.
\item Dire si le gestionnaire pourrait vendre des parts d'ETF à la place de l'obligation.
\item Expliquer ce qu'un \omterm{def:m2:how-bonds-trade:composite}{prix composite} apprend au gestionnaire avant qu'il ne demande une cotation.
\end{enumerate}

\medskip\noindent\textbf{Partie IV --- Jugement.}
\begin{enumerate}[resume]
\item Expliquer pourquoi les dealers n'affichent pas toujours leur meilleur cours à chaque client.
\item Expliquer comment mesurer la fuite dans ses propres registres de transactions.
\item Expliquer pourquoi le meilleur nombre pourrait différer pour une obligation high yield.
\item Énoncer le \emph{résultat nommé}~: le nombre de dealers qui minimise le coût
espéré, ainsi que ce coût.
\item Formuler en une phrase pourquoi interroger davantage de dealers n'est pas toujours préférable.
\end{enumerate}
\end{problem}

\section{Questions d'entretien}

\begin{interviewq}[$\star$ \iqroles{trader, développeur}]\label{iq:m2:how-bonds-trade:1}
Comment une obligation d'entreprise s'achète-t-elle et se vend-elle, par comparaison avec une action~?
\end{interviewq}

\begin{interviewq}[$\star$ \iqroles{trader, chercheur}]\label{iq:m2:how-bonds-trade:2}
Qu'est-ce que TRACE, et que ne montre-t-elle pas~?
\end{interviewq}

\begin{interviewq}[$\star\star$ \iqroles{chercheur}]\label{iq:m2:how-bonds-trade:3}
Combien de dealers un client devrait-il interroger dans une RFQ~? Construisez un modèle simple.
\end{interviewq}

\begin{interviewq}[$\star\star$ \iqroles{trader, banque}]\label{iq:m2:how-bonds-trade:4}
Comment un dealer fixe-t-il le prix d'une \omterm{def:m2:how-bonds-trade:a2a}{transaction de portefeuille} de 400 obligations~?
\end{interviewq}

\begin{interviewq}[$\star\star$ \iqroles{chercheur, trader}]\label{iq:m2:how-bonds-trade:5}
Pourquoi un ETF obligataire peut-il se négocier avec une décote par rapport à sa valeur liquidative, et
quel prix est le bon~?
\end{interviewq}

\begin{interviewq}[$\star\star\star$ \iqroles{développeur}]\label{iq:m2:how-bonds-trade:6}
Concevez un système qui répond à 20\,000 RFQ obligataires par jour avec des prix dans un délai de 100
millisecondes.
\end{interviewq}
